\usepackage[paperwidth=176mm,paperheight=250mm,left=22mm,right=22mm,top=22mm,bottom=23mm,headheight=15pt,headsep=7mm]{geometry}
\usepackage{fontspec}
\usepackage{xeCJKfntef}
\IfFontExistsTF{SimSun}{\setCJKmainfont{SimSun}[AutoFakeBold=2]\setCJKsansfont{SimHei}\setCJKmonofont{FangSong}\newCJKfontfamily\quotefont{FangSong}\newCJKfontfamily\heiti{SimHei}}{\setCJKmainfont{FandolSong-Regular}[BoldFont=FandolSong-Bold]\setCJKsansfont{FandolHei-Regular}\setCJKmonofont{FandolFang-Regular}\newCJKfontfamily\quotefont{FandolFang-Regular}\newCJKfontfamily\heiti{FandolHei-Regular}}
\setmainfont{TeX Gyre Termes}
\IfFontExistsTF{Calibri}{\setsansfont{Calibri}}{\setsansfont{TeX Gyre Heros}}
\usepackage{graphicx,caption,subcaption,booktabs,longtable,array,enumitem,fancyhdr,needspace}
\usepackage{xurl}
\usepackage{amsthm}
\usepackage{tcolorbox}
\tcbuselibrary{skins,breakable}
% 读者注释、批注与评论；可选参数用于局部调整 tcolorbox 样式。
% tcolorbox 依赖的 verbatim 已定义同名隐藏环境，此处改为可见评论框。
\renewtcolorbox{comment}[1][]{
  enhanced,breakable,frame hidden,colback=white,
  boxrule=.35pt,arc=2mm,
  borderline={.35pt}{0pt}{black!60,dash pattern=on 2pt off 2pt},
  left=3mm,right=3mm,top=2mm,bottom=2mm,
  before skip=\BookParSkip,after skip=\BookParSkip,
  fontupper=\small\quotefont,
  before upper={\setlength{\parindent}{0pt}\setlength{\parskip}{\BookParSkip}},
  #1
}
\usepackage[unicode,hidelinks,bookmarksnumbered]{hyperref}
\hypersetup{pdftitle={真事隐：康熙废储与正史虚构},pdfauthor={孙立天},pdfsubject={扫描书 B5 重排}}
\ctexset{chapter={format=\raggedright\sffamily\heiti\zihao{2},beforeskip=20pt,afterskip=28pt},section={format=\raggedright\sffamily\heiti\zihao{4},numbering=true,beforeskip=22pt,afterskip=14pt}}
\setcounter{secnumdepth}{1}
\setcounter{tocdepth}{1}
\linespread{1.25}
\setlength{\parindent}{2em}
\newlength{\BookParSkip}
\setlength{\BookParSkip}{.35\baselineskip}
\setlength{\parskip}{\BookParSkip}
% 定律按章编号；标题黑体及无衬线西文，正文保持宋体正体。
\newtheoremstyle{booklaw}
  {\BookParSkip}{\BookParSkip}
  {\normalfont}{0pt}{\sffamily\heiti}{。}{.5em}
  {\thmname{#1}\thmnumber{ #2}\thmnote{（#3）}}
\theoremstyle{booklaw}
\newtheorem{law}{定律}[chapter]
\newtheorem*{law*}{定律}
\emergencystretch=2em
% 实见字符暂不可辨时的可见占位；框宽分别为一个汉字与半个西文字身。
\newcommand{\UnknownGlyphBox}[1]{\begingroup
  \setlength{\fboxsep}{0pt}\setlength{\fboxrule}{.4pt}%
  \fbox{\makebox[\dimexpr#1-2\fboxrule\relax]{%
    \rule[-.15em]{0pt}{\dimexpr1em-2\fboxrule\relax}}}%
\endgroup}
\DeclareRobustCommand{\UnclearCJK}{\UnknownGlyphBox{1em}}
\DeclareRobustCommand{\UnclearLatin}{\UnknownGlyphBox{.5em}}
\newenvironment{bookquote}{\list{}{\leftmargin=2em\rightmargin=2em\listparindent=2em\itemindent=2em\parsep=\BookParSkip\topsep=0pt\partopsep=0pt}\item\relax\quotefont}{\endlist}
\makeatletter
\newcommand{\NotesLocation}{\if@mainmatter\ifnum\value{section}>0 \thesection\else 第\chinese{chapter}章\fi\else\leftmark\fi}
\makeatother
\newenvironment{sectionnotes}{\par\needspace{4\baselineskip}\medskip\noindent{\sffamily\heiti\zihao{4} 注释（\NotesLocation）}\par\smallskip\begingroup\small\setlength{\parindent}{0pt}\setlength{\parskip}{.7\baselineskip}\linespread{1.12}\selectfont}{\par\endgroup\medskip}
\newcommand{\NoteRef}[2]{\hyperlink{note:#1:#2}{\textsuperscript{#2}}}
\newcommand{\BookNote}[3]{\par\hypertarget{note:#1:#2}{}\hangindent=2.8em\hangafter=1\noindent\makebox[2.8em][l]{#2}#3\par}
\newcommand{\BibEntry}[1]{\par\hangindent=2em\hangafter=1\noindent #1\par\smallskip}
\captionsetup{font=small,labelfont=bf,labelsep=quad}
\captionsetup[subfigure]{font=small,labelformat=simple}
\renewcommand{\thesubfigure}{\arabic{subfigure}}
\renewcommand{\figurename}{图片}
\renewcommand{\thefigure}{\arabic{figure}}
\usepackage{chngcntr}
\counterwithout{figure}{chapter}
\renewcommand{\tablename}{表}
\renewcommand{\thetable}{\arabic{table}}
\counterwithout{table}{chapter}
\pagestyle{fancy}\fancyhf{}
\fancyfoot[L]{\small\quotefont\nouppercase{\leftmark}}
\fancyfoot[R]{\thepage}
\renewcommand{\headrulewidth}{0pt}
\fancypagestyle{plain}{\fancyhf{}\fancyfoot[L]{\small\quotefont\nouppercase{\leftmark}}\fancyfoot[R]{\thepage}\renewcommand{\headrulewidth}{0pt}}
\newcommand{\ScanPage}[1]{\clearpage\newgeometry{margin=0pt}\thispagestyle{empty}\noindent\includegraphics[width=\paperwidth,height=\dimexpr\paperheight-1pt\relax,keepaspectratio]{#1}\par\clearpage\restoregeometry}
